\documentclass[12pt]{article}
\usepackage[utf8]{inputenc}
\usepackage[T1]{fontenc}
\usepackage{hyperref}
\title{Flaaitaal Language Resources for K-12 Curriculum Development}
\author{Compiled by Claude Sonnet 5.5 (claude-sonnet-5-5) for Andr\'as Kornai}
\date{}
\begin{document}
\maketitle

\section{Introduction}
Flaaitaal (also Tsotsitaal, Fly Taal, ISO 639-3 \texttt{fly}) is an urban speech register of South Africa with an Afrikaans-based lexical core. It has no ethnic community of its own and is used as a second language. Searches found no school materials; what exists is mostly dictionaries and academic description. Sources disagree on naming, and some speakers consider ``tsotsitaal'' outdated or insulting.

\section{Primary Educational Resources}
No primers, textbooks or school programmes for Flaaitaal were found in search results.

\section{Linguistic and Reference Materials}
\begin{itemize}
\item Dictionary of South African English entry for \emph{flaaitaal}: \url{https://dsae.co.za/entry/flaaitaal/e02397}
\item Review in \emph{Lexikos} of Molamu's \emph{Tsotsitaal: A Dictionary of the Language of Sophiatown} (UNISA, 2003); the review notes there is no formal orthography: \url{https://lexikos.journals.ac.za/pub/article/view/743/339}
\item Glottolog reference record for Makhudu (1995), ``An Introduction to Flaaitaal'': \url{https://glottolog.org/resource/reference/id/550543}
\item Cambridge book chapter on Flaaitaal/Tsotsitaal (Makhudu and Mesthrie): \url{https://www.cambridge.org/core/product/identifier/CBO9780511486692A032/type/BOOK_PART}
\end{itemize}

\section{Religious and Community Resources}
Joshua Project reports no Bible translations or recordings: \url{https://joshuaproject.net/languages/fly}

\section{Digital and Audio Resources}
Glosbe hosts a community-contributed Flaaitaal--English dictionary of uneven, unverified quality: \url{https://glosbe.com/fly/en}

\section{Curriculum Development Context}
Little exists. Any heritage or revitalization curriculum would have to be built from scratch, probably as a language-awareness unit on urban registers, using the dictionary and academic sources above with care, since no standard orthography exists.

\end{document}